\newcommand{\GeneralInteractionRecoilFigure}{%
\begin{figure}[htbp]
\centering
\begin{tikzfig}[x=1cm,y=1cm,>=Stealth]
\node[lecturebox,text width=3.15cm,minimum height=1.35cm] (QED) at (-4.2,1.55)
{微观电磁耦合\\正能电子投影};
\node[lecturebox,text width=3.15cm,minimum height=1.35cm] (CONT) at (0,1.55)
{连续电子与环境\\$|\vb{p},s\rangle\otimes|\mu\rangle$};
\node[lecturebox,text width=3.15cm,minimum height=1.35cm] (SEL) at (4.2,1.55)
{有限时间跃迁\\矩阵元与失谐};

\node[lecturebox,text width=3.15cm,minimum height=1.35cm] (REC) at (4.2,-1.2)
{保留反冲的通道\\使用真实能量 $E(\vb{p}_\ell)$};
\node[lecturebox,text width=3.15cm,minimum height=1.35cm] (UNI) at (0,-1.2)
{等间隔能量格\\受控的低反冲极限};

\draw[->] (QED.east)--(CONT.west);
\draw[->] (CONT.east)--(SEL.west);
\draw[->] (SEL.south)--(REC.north);
\draw[->] (REC.west)--(UNI.east);
\end{tikzfig}
\caption{从连续电子态到离散交换通道。电子电流顶角与环境谱确定耦合，有限作用时间和能量匹配决定哪些通道有效参与交换。保留各通道的真实能量 $E(\vb{p}_\ell)$，得到反冲格；当各条边的失谐和耦合差别都可忽略时，才得到等间隔能量格。}
\label{fig:general_interaction_continuum_to_recoil_ladder}
\end{figure}}
